% Generated draft: requires mathematical and source review.
\documentclass[11pt]{article}
\usepackage[margin=25mm]{geometry}
\usepackage{fontspec}
\setmainfont{Latin Modern Roman}
\usepackage{amsmath,amssymb,mathtools,unicode-math}
\setmathfont{Latin Modern Math}
\usepackage{enumitem,needspace,xurl,hyperref}
\setlength{\parindent}{0pt}
\setlength{\parskip}{5pt}
\setlength{\emergencystretch}{3em}
\newcommand{\N}{\mathbb{N}}
\newcommand{\Z}{\mathbb{Z}}
\newcommand{\Q}{\mathbb{Q}}
\newcommand{\R}{\mathbb{R}}
\newcommand{\C}{\mathbb{C}}
\newcommand{\F}{\mathbb{F}}
\newcommand{\A}{\mathbb{A}}
\newcommand{\PP}{\mathbb{P}}
\newcommand{\E}{\mathbb{E}}
\newcommand{\eps}{\varepsilon}
\newcommand{\dd}{\,\mathrm{d}}
\begin{document}
\section*{UnsolvedMath ID: 801182}
\textbf{Source problem code:} catalogue entry
\subsection{Definitions and mathematical statement}
Let \(\Lambda_{\Q}\) be the graded ring of symmetric functions over \(\Q\). A partition \(\lambda\vdash d\) is a finite weakly decreasing sequence of positive integers whose sum is \(d\); the empty partition has size zero. Write \(s_\lambda\) for the corresponding Schur function, and equip \(\Lambda_{\Q}\) with the Hall inner product, normalized by
\[
\langle s_\lambda,s_\nu\rangle=\delta_{\lambda,\nu}.
\]
For \(r\geq1\), let \(p_r=\sum_{i\geq1}x_i^r\). For \(f,g\in\Lambda_{\Q}\), plethysm \(f[g]\) is obtained by expressing \(f\) in the power sums and replacing each \(p_r\) by the symmetric function obtained from \(g\) under the substitution \(p_k\mapsto p_{rk}\), with rational coefficients fixed. For a positive integer \(m\), write \(s_m=s_{(m)}\).

The explicit target in the supplied record is the following problem, attributed to Stanley (1999), Problem 9. For every pair of positive integers \(m,n\) and every partition \(\lambda\vdash mn\), give a positive combinatorial interpretation of
\[
a_{m,n}^{\lambda}:=\langle s_m[s_n],s_\lambda\rangle,
\qquad
s_m[s_n]=\sum_{\lambda\vdash mn}a_{m,n}^{\lambda}s_\lambda.
\]
More precisely, describe finite sets \(\mathcal C_{m,n}^{\lambda}\) by explicit combinatorial conditions and establish
\[
a_{m,n}^{\lambda}=\#\mathcal C_{m,n}^{\lambda},
\]
so that the interpretation itself supplies a combinatorial proof of the stated nonnegativity. Homogeneity gives coefficient zero for partitions whose size differs from \(mn\).

\textit{Unresolved qualifications.} The source does not specify the range of \(m,n\); the formulation above makes explicit the usual positive-integer reading. Whether zero indices are intended remains unspecified. The source also mentions general plethysms \(s_\alpha[s_\beta]\), whose coefficients are \(\langle s_\alpha[s_\beta],s_\nu\rangle\) for \(\nu\vdash |\alpha||\beta|\), but its explicit request concerns the single-row case \(s_m[s_n]\). It is unclear whether a general positive rule is additionally required. The supplied record labels the problem open; this draft does not independently verify that status.
\subsection{Short English statement}
For positive integers \(m,n\), find explicitly described finite combinatorial objects counted by each Schur coefficient of \(s_m[s_n]\). Such a counting rule should itself explain the supplied assertion that these coefficients are nonnegative. The source leaves zero-index conventions and the intended scope of its mention of general Schur plethysms unresolved.
\subsection{Sources}
\begin{itemize}
\item Supplied local UnsolvedMath registry record: permanent ID \(801182\), exact source problem code \texttt{catalogue entry}, in the AMR Open Problem Lists collection. The record attributes the problem to Stanley, with proposed year 1999, and labels its status \texttt{open}.
\item Stanley, \emph{Positivity Problems and Conjectures in Algebraic Combinatorics} (1999), Problem 9, as identified by the supplied record: \url{https://math.mit.edu/~rstan/papers/problems.pdf}. The supplied provenance records access on 2026-10-16 and extraction method \texttt{pdf+html\_verified}; these are supplied provenance claims, not independent source verification by this draft.
\item The supplied background describes this as one of the central open problems in algebraic combinatorics. This description and the open-status label are retained as source assessments, not independently verified literature conclusions. The registry record was updated on 2026-10-16.
\end{itemize}
\end{document}
